\subsection{CHARACTERS OF g-MODULES}

Let \(\Delta\) be a commutative monoid written additively, and \(\mathbf{Z}[\Delta] = \mathbf{Z}^{(\Delta)}\) the algebra of the monoid \(\Delta\) over \(\mathbf{Z}\) (Algebra, Chap. III, §2, no. 6). Denote by \((e^{\lambda})_{\lambda \in \Delta}\) the canonical basis of \(\mathbf{Z}[\Delta]\) . For all \(\lambda, \mu \in \Delta\) , we have \(e^{\lambda + \mu} = e^{\lambda}e^{\mu}\) . If 0 is the neutral element of \(\Delta\) , then \(e^0\) is the unit element of \(\mathbf{Z}[\Delta]\) ; it is denoted by 1.

Let E be a \(\Delta\) -graded vector space over a field \(\kappa\) , and let \((\mathrm{E}^{\lambda})_{\lambda\in\Delta}\) be its graduation. If each \(E^{\lambda}\) is finite dimensional, the character of E, denoted by \(\operatorname{ch}(\operatorname{E})\) , is the element \((\dim\operatorname{E}^{\lambda})_{\lambda\in\Delta}\) of \(Z^{\Delta}\) . If E itself is finite dimensional,

\[
\operatorname{ch} (\mathrm{E}) = \sum_ {\lambda \in \Delta} (\dim \mathrm{E} ^ {\lambda}) e ^ {\lambda} \in \mathbf {Z} [ \Delta ].\tag{5}
\]

Let \(E', E, E''\) be \(\Delta\) -graded vector spaces such that the \(E'^{\lambda}, E^{\lambda}, E''^{\lambda}\) are finite dimensional over \(\kappa\) , and \(0 \longrightarrow E' \longrightarrow E \longrightarrow E'' \longrightarrow 0\) an exact sequence of graded homomorphisms of degree 0. It is immediate that

\[
\operatorname{ch} (\mathrm{E}) = \operatorname{ch} (\mathrm{E} ^ {\prime}) + \operatorname{ch} (\mathrm{E} ^ {\prime \prime}).\tag{6}
\]

In particular, if \(F_{1}, F_{2}\) are \(\Delta\) -graded vector spaces such that the \(F_{1}^{\lambda}\) and the \(F_{2}^{\lambda}\) are finite dimensional over \(\kappa\) , then

\[
\operatorname{ch} \left(\mathrm{F} _ {1} \oplus \mathrm{F} _ {2}\right) = \operatorname{ch} \left(\mathrm{F} _ {1}\right) + \operatorname{ch} \left(\mathrm{F} _ {2}\right).\tag{7}
\]

If \(F_{1}\) and \(F_{2}\) are finite dimensional, we also have

\[
\operatorname{ch} \left(\mathrm{F} _ {1} \otimes \mathrm{F} _ {2}\right) = \operatorname{ch} \left(\mathrm{F} _ {1}\right). \operatorname{ch} \left(\mathrm{F} _ {2}\right).\tag{8}
\]

Example. Assume that \(\Delta = N\) . Let T be an indeterminate. There exists a unique isomorphism from the algebra Z{[}N{]} to the algebra Z{[}T{]} that takes \(e^{n}\) to \(T^{n}\) for all \(n \in N\) . For any finite dimensional N-graded vector space E, the image of \(\operatorname{ch}(E)\) in Z{[}T{]} is the Poincaré polynomial of E (Chap. V, §5, no. 1).

Let E be a g-module such that \(E = \sum_{\lambda \in h^{*}} E^{\lambda}\) and such that each \(E^{\lambda}\) is finite dimensional. We know that \((\mathrm{E}^{\lambda})_{\lambda \in \mathfrak{h}^{*}}\) is a graduation of the vector space E. In what follows we shall reserve the notation \(\operatorname{ch}(\mathrm{E})\) for the character of E considered as a \(h^{*}\) -graded vector space. Thus, the character \(\operatorname{ch}(\mathrm{E})\) is an element of \(Z^{h^{*}}\) . If E is finite dimensional, \(\operatorname{ch}(\mathrm{E}) \in \mathbf{Z}[\mathrm{P}]\) . By formula (6) and Lemma 3 of no. 6, there exists a unique homomorphism from the group \(\mathcal{R}(\mathfrak{g})\) to Z{[}P{]} that takes E to \(\operatorname{ch}(\mathrm{E})\) , for any finite dimensional g-module E; this homomorphism will be denoted by ch.~Relation (8) shows that ch is a homomorphism from the ring \(\mathcal{R}(\mathfrak{g})\) to the ring Z{[}P{]}.

Remark. Every element of P defines a simple h-module of dimension 1, hence a homomorphism from the group Z{[}P{]} to the group \(\mathcal{R}(\mathfrak{h})\) , which is an injective homomorphism of rings. It is immediate that the composite

\[
\mathcal {R} (\mathfrak {g}) \longrightarrow \mathbf {Z} [ \mathrm{P} ] \longrightarrow \mathcal {R} (\mathfrak {h})
\]

is the homomorphism defined by the canonical injection of \(\mathfrak{h}\) into \(\mathfrak{g}\) (no. 6).

The Weyl group W operates by automorphisms on the group P, and hence operates on \(Z^{P}\) . For all \(\lambda \in P\) and all \(w \in W\) , we have \(we^{\lambda} = e^{w\lambda}\) . Let \(Z[P]^{W}\) be the subring of Z{[}P{]} consisting of the elements invariant under W.

Lemma 6. If \(\lambda \in \mathrm{P}_{++}\) , then \(\operatorname{ch}[\lambda] \in \mathbf{Z}[\mathrm{P}]^{\mathrm{W}}\) . The unique maximal term of \(\operatorname{ch}[\lambda]\) (Chap. VI, §3, no. 2, Def. 1) is \(e^{\lambda}\) .

The first assertion follows from no. 1, Cor. 2 of Prop. 2, and the second from §6, no. 1, Prop. 1 (ii).

THEOREM 2. (i) Let \((\varpi_{\alpha})_{\alpha \in \mathrm{B}}\) be the family of fundamental weights relative to B. Let \((\mathrm{T}_{\alpha})_{\alpha \in \mathrm{B}}\) be a family of indeterminates. The map \(f\mapsto f(\left([ \varpi_{\alpha} ]\right)_{\alpha \in \mathrm{B}})\) from \(\mathbf{Z}[(\mathrm{T}_{\alpha})_{\alpha \in \mathrm{B}}]\) to \(\mathcal{R}(\mathfrak{g})\) is an isomorphism of rings.

\begin{enumerate}
\def\labelenumi{(\roman{enumi})}
\setcounter{enumi}{1}
\item
  The homomorphism \(\operatorname{ch}\) from \(\mathcal{R}(\mathfrak{g})\) to \(\mathbf{Z}[\mathrm{P}]\) induces an isomorphism from the ring \(\mathcal{R}(\mathfrak{g})\) to the ring \(\mathbf{Z}[\mathrm{P}]^{\mathrm{W}}\) .
\item
  Let E be a finite dimensional g-module. If \(\operatorname{ch} E = \sum_{\lambda \in P_{++}} m_{\lambda} \operatorname{ch}[\lambda]\) , the isotypical component of E of highest weight \(\lambda\) has length \(m_{\lambda}\) .
\end{enumerate}

The family \(([\lambda])_{\lambda \in \mathrm{P}_{++}}\) is a basis of the \(\mathbf{Z}\) -module \(\mathcal{R}(\mathfrak{g})\) , and the family \((\operatorname{ch}[\lambda])_{\lambda \in \mathrm{P}_{++}}\) is a basis of the \(\mathbf{Z}\) -module \(\mathbf{Z}[\mathrm{P}]^{\mathrm{W}}\) (Lemma 6, and Chap. VI, §3, no. 4, Prop. 3). This proves (ii) and (iii). Assertion (i) follows from (ii), Lemma 6 and Chap. VI, §3, no. 4, Th. 1.

COROLLARY. Let E, E' be finite dimensional g-modules. Then E is isomorphic to E' if and only if ch E = ch E'.

This follows from Th. 2 (ii) and the fact that E, E' are semi-simple.

